%%%%%%%%%%%%%%%%%%%%%%%%%%%%%%%%%
%
%       EXERCÍCIO
%
%%%%%%%%%%%%%%%%%%%%%%%%%%%%%%%%%

\definevalorquestao

\begin{exercicioBanco}[\valorquestao]
Uma biblioteca escolar registrou a quantidade de livros emprestados nos últimos cinco meses:

\begin{itemize}
    \begin{multicols}{3}
        \item 1º mês: 18;
        \item 2º mês: 20;
        \item 3º mês: 24;
        \item 4º mês: 26;
        \item 5º mês: 22.
    \end{multicols}
\end{itemize}

No início do primeiro mês, a biblioteca possuía 210 livros disponíveis para empréstimo. A política de reposição prevê a aquisição de novos livros, no início do sexto mês, de modo que a quantidade inicial disponível para os próximos meses seja igual a 12 vezes a média das quantidades mensais emprestadas nos últimos cinco meses. Qual é a quantidade de livros que deve ser adquirida no início do sexto mês?
\end{exercicioBanco}

